9.3 Compact Sets
15
b) Not every $t_1$-space is a $t_2$-space. (See Example 4.)
c) The two-point space $X = \{a, b\}$ with the open sets $\{\emptyset, X\}$ is not a $t_1$-space.
d) In a $t_1$-space a set $F$ is closed if and only if it contains all its limit points.
7. a) Prove that in any topological space there is an everywhere dense set whose
cardinality does not exceed the weight of the space.
b) Verify that the following metric spaces are separable: $C[a, b]$, $C^{(k)}[a, b]$,
$R_1[a, b]$, $R_p[a, b]$ (for the formulas giving the respective metrics see Sect. 9.1).
c) Verify that if $\max$ is replaced by $\sup$ in relation (9.6) of Sect. 9.1 and regarded
as a metric on the set of all bounded real-valued functions defined on a closed inter-
val $[a, b]$, we obtain a nonseparable metric space.

9.3 Compact Sets

9.3.1 Definition and General Properties of Compact Sets
Definition 1 A set $K$ in a topological space $(X, \tau)$ is compact (or bicompact$^3$) if
from every covering of $K$ by sets that are open in $X$ one can select a finite number
of sets that cover $K$.
Example 1 An interval $[a, b]$ of the set $\mathbb{R}$ of real numbers in the standard topology
is a compact set, as follows immediately from the lemma of Sect. 2.1.3 asserting
that one can select a finite covering from any covering of a closed interval by open
intervals.
In general an $m$-dimensional closed interval $I^m = \{x \in \mathbb{R}^m \mid a^i \leq x^i \leq b^i, i = 1,
\ldots, m\}$ in $\mathbb{R}^m$ is a compact set, as was established in Sect. 7.1.3.
It was also proved in Sect. 7.1.3 that a subset of $\mathbb{R}^m$ is compact if and only if it
is closed and bounded.
In contrast to the relative properties of being open and closed, the property of
compactness is absolute, in the sense that it is independent of the ambient space.
More precisely, the following proposition holds.
Proposition 1 A subset $K$ of a topological space $(X, \tau)$ is a compact subset of $X$ if
and only if $K$ is compact as a subset of itself with the topology induced from $(X, \tau)$.
Proof This proposition follows from the definition of compactness and the fact that
every set $G_K$ that is open in $K$ can be obtained as the intersection of $K$ with some
set $G_X$ that is open in $X$.
$^3$The concept of compactness introduced by Definition 1 is sometimes called bicompactness in
topology.